\documentclass{amsart}
\input{C:/Users/jzchr/MathDocuments/preamble_general.tex}

\renewcommand{\X}{\mathfrak{X}}

\title{Math 319 HW 1}
\author{Jalen Chrysos}
\begin{document}
	\maketitle 
	
	\textbf{Problem 2}: Consider the vector fields
	$$X_1(p) := e_1 - \<e_1,p\>p, \;\;\; X_2(p) := e_2 - \<e_2,p\>p$$
	on $S^3$. Compute $\nabla_{X_2}X_1(q)$ where $q=(\tfrac12,\tfrac12,\tfrac12,\tfrac12)$
	\begin{proof}
		We saw that the Levi-Civita connection on $S^n$ is the orthogonal projection of the one on $\R^{n+1}$, i.e.
		$$
		\nabla^{S^n}_{X_2}X_1 = (\nabla^{\R^{n+1}}_{X_2^{ext}}X_1^{ext})^T.
		$$
		First we compute $\nabla^{\R^4}_{X_2}X_1$ at $p$. Letting $\partial_i = \partial/\partial e_i$, because Christoffel symbols are all 0 in the Euclidean metric, we have
		$$
		\nabla^{\R^4}_{X_2}X_1 = \sum_k X_2(a_k) \partial_k
		$$
		where $X_1 = \sum_k a_k \partial_k$. More specifically, $X_1$ is
		$$
		X_1(p) = (1 - p_1^2,-p_1p_2, -p_1p_3, -p_1p_4)
		$$ 
		and so
		$$
		\partial_{p_1} X_1 = (-2p_1,-p_2,-p_3,-p_4), \; \partial_{p_2} X_1 = (0,-p_1,0,0), \; \partial_{p_3} X_1 = (0,0,-p_1,0),\;  \partial_{p_4} X_1 = (0,0,0,-p_1).
		$$
		At $q$, these are
		$$
		(-1,-\tfrac12,-\tfrac12,-\tfrac12), (0,-\tfrac12,0,0),(0,0,-\tfrac12,0),(0,0,0,-\tfrac12).
		$$
		$X_2(q)=(-\tfrac14,\tfrac34,-\tfrac14,-\tfrac14)$, so we have
		\begin{align*}
		(\nabla_{X_2}^{\R^4} X_1)(q) &= (X_2(a_1),X_2(a_2),X_2(a_3),X_2(a_4))\\
		 &=  -\tfrac14 (-1,-\tfrac12,-\tfrac12,-\tfrac12) + \tfrac34(0,-\tfrac12,0,0) + -\tfrac14(0,0,-\tfrac12,0) + -\tfrac14(0,0,0,-\tfrac12)\\
		 &= (\tfrac14,-\tfrac14,\tfrac14,\tfrac14).
		\end{align*}
		Now projecting onto $T_qS^3$, we get
		$$
		\nabla_{X_2}^{S^3}(X_1) = (\tfrac14,-\tfrac14,\tfrac14,\tfrac14)^T = (\tfrac14,-\tfrac14,\tfrac14,\tfrac14) - \<(\tfrac14,-\tfrac14,\tfrac14,\tfrac14),q\>q = (\tfrac18,-\tfrac38,\tfrac18,\tfrac18).
		$$
	\end{proof}\\
	
	
	\newpage 
	
	\textbf{Problem 7 (do Carmo 2.5)}: In Euclidean space, the parallel transport of a vector between two points does not depend on the curve joining the two points. Show that this is not necessarily true on an arbitrary Riemannian Manifold.
	\begin{proof}
		Let $M=S^2$, and consider the two points $p = (0,0,1)$ and $q=(0,0,-1)$, the north and south poles of $S^2$. Let $\vec{v} =(1,0,0) \in T_pS^2$. We will transport this tangent vector from $p$ to $q$ along two curves: 
		$$
		\gamma_1:t\mapsto (0,\sin t,\cos t), \;\;\; \gamma_2:t\mapsto (\sin t,0,\cos t).
		$$
		$\gamma_1(0)=\gamma_2(0)=p$ and $\gamma_1(\pi)=\gamma_2(\pi)=q$, so both are paths from $p$ to $q$. Consider the following vector fields, $X_1$ along $\gamma_1$ and $X_2$ along $\gamma_2$:
		$$
		X_1:t\mapsto (1,0,0), \;\;\; X_2:t\mapsto (\cos t, 0 ,-\sin t).
		$$
		These are the parallel transports of $(1,0,0)$ along $\gamma_1,\gamma_2$: first, $X_1$ is constant, so
		$$
		DX_1/\d t = \nabla^{S^2}_{\gamma_1'} X_1 = (\nabla^{\R^3}_{\gamma_1'} X_1)^T = 0^T = 0.
		$$
		And $X_2(t)$ is the tangent vector field along $\gamma_2$, which is a geodesic, so $DX_2/\d t=0$ as well.
%		\begin{align*}
%		\frac{DX_2}{\d t} &= \nabla_{\d \gamma_2/\d t} X_2\\
%		&= \nabla_{X_2} X_2 \\
%		&= 0
%		\end{align*}
%		$\nabla_AA=0$ follows from the defining property of the Levi-Civita connection:
%		\begin{align*}
%		\<Z,\nabla_A A\> &= \frac12 \Big(A\<A,Z\> + A\<Z,A\> - Z\<A,A\> - \<[A,Z],A\> - \<[A,Z],A\> - \<[A,A],Z\>\Big)\\
%		&= \frac12 \Big(A\<A,Z\> + A\<Z,A\> - Z|A|^2 - 2\<[A,Z],A\>\Big)\\
%		\end{align*}
		But $X_1(\pi)=(1,0,0)$ while $X_2(\pi)=(-1,0,0)$, so parallel transport can result in different outcomes depending on the curve taken.
	\end{proof}\\
	
	\newpage
	
	\textbf{Problem 9 (do Carmo 3.5(d))}: Let $M$ be a Riemannian manifold and $X\in \X(M)$. Let $p\in M$ and $U\subset M$ be a neighborhood of $p$. Let $\ph:(-\ep,\ep)\times U\to M$ be the local flow generated by $X$. $X$ is called a \textit{Killing field} if for $t'\in (-\ep,\ep)$, the map $\ph(t',\bullet):U\to M$ is an isometry. Show that $X$ is a Killing field iff 
	$$
	\<\nabla_YX, Z\> + \<\nabla_Z X,Y\>=0
	$$
	for all $Y,Z\in \X(M)$.
	\begin{proof}
		Using the fact that $\nabla$ is the Levi-Civita connection, we have the formula
		$$
		\<\nabla_YX,Z\> = \frac12 \Big(Y\<X,Z\> + X\<Z,Y\> - Z\<X,Y\> - \<[X,Y],Z\> - \<[X,Z],Y\> - \<[Y,Z],X\>\Big)
		$$
		so we can put $\<\nabla_Y X,Z\>$ and $\<\nabla_Z X,Y\>$ in this form and add them, yielding
		$$
		\<\nabla_Y X,Z\> + \<\nabla_Z X,Y\> = X\<Y,Z\> - \<[X,Y],Z\> - \<[X,Z],Y\>.
		$$
		Now let's assume that $X=\partial_t$, $Y=\partial_i$, and $Z=\partial_j$ in local coordinates, so that $X$ is orthogonal to both $Y$ and $Z$. In this case $[\partial_t,\partial_i]=[\partial_t,\partial_j]=0$, so we have
		$$
		\<\nabla_Y X,Z\> + \<\nabla_Z X, Y\> = X\<Y,Z\> = \partial_t \<\partial_i,\partial_j\>.
		$$
		Because $X$ is a Killing field, the flow along $X$ preserves the metric, i.e. $\partial_t \<\partial_i,\partial_j\>=0$. So
		$$
		\<\nabla_{\partial_i} X,\partial_j \> + \<\nabla_{\partial_j} X, \partial_i\> = 0 
		$$
		as expected. To extend this to all $Y$ and $Z$, letting $Y = \sum_i y_i \partial_i$ and $Z=\sum_j z_j\partial_j$, we have
		$$
		\<\nabla_Y X,Z\> + \<\nabla_ZX,Y\> = \sum_{ij} y_iz_j\Big(\<\nabla_{\partial_i} X, \partial_j\> + \<\nabla_{\partial_j} X, \partial_i\> \Big) = 0.
		$$
		
		In the converse direction, we similarly see that if the above expression is 0 for all $y_i,z_j\in C^{\infty}(M)$, then $\<\nabla_{\partial_i} X,\partial_j\> + \<\nabla_{\partial_j} X,\partial_i\>=0$ for all $i,j$. And again, this is equal to $X\<\partial_i,\partial_j\>$, so if it is 0 then the flow along $X$ is an isometry, i.e. $X$ is a Killing field.
	\end{proof}\\
	
	\newpage
	\textbf{Problem 10 (do Carmo 3.6)}: Let $X$ be a Killing field on a connected Riemannian manifold $M$. Assume that there is a point $q\in M$ such that $X(q)=0$ and $\nabla_YX(q)=0$ for all $Y$. Show that $X=0$.
	\begin{proof}
		Let $\ph_t$ be the flow associated with $X$. We are given that $X(q)=0$, so $\ph_t(q)=q$ for all $t$. This implies $\nabla_X Y(q)=0$. And we are also given that $\nabla_Y X(q)=0$ for all $Y$, so in fact
		$$
		[X,Y](q) = (\nabla_X Y - \nabla_Y X) = 0-0 = 0.
		$$
		Now, writing $[X,Y]$ in terms of the flow definition,
		$$
		[X,Y](q) = \partial_t|_{t=0} \d\ph_{-t} Y_{\ph_t(q)} = \partial_t|_{t=0} \d\ph_{-t} Y_{q}
		$$
		As this is 0 for all $Y_q$, we see that $\d\ph_{-t}$ and hence $\d\ph_t$ is constant in $t$ on all vectors $Y_q$. This also implies that $\d\ph_t$ is the identity, since $\d\ph_t(Y_q) = \d\ph_{t+s}(Y_q)=\d\ph_t\cdot \d\ph_s(Y_q)$ for all $s\in \R$, hence $\d\ph_s(Y_q)=Y_q$.\\
		
		Now, since we know $\d\ph_t = I$, $\ph_t(\exp_q)=\exp_q(\d\ph_t) = \exp_q$, so $\ph_t$ acts as the identity on the image of $\exp_q$. By the normal neighborhood theorem, the image of $\exp_q$ contains an open neighborhood around $q$, so $X$ is 0 on an open neighborhood $U$ of $q$. Now because $M$ is connected, every $x\in M$ is connected to $q$ by a (possibly non-unique) shortest path, which is a geodesic $\gamma$. $\ph_t$ acts as the identity on part of this geodesic (the part in $U$) hence it must preserve the whole geodesic, because $\ph_t$ is an isometry and maps geodesics to geodesics, and $\gamma$ is determined by its behavior in $U$ which is fixed by $\ph_t$, so $\ph_t(\gamma)=\gamma$. Thus $\ph_t(x)=x$ as well, so $\ph_t$ is the identity everywhere, and hence $X=0$.
	\end{proof}
\end{document}